\section{ERI Task and Benchmark}
\label{sec:task-benchmark}

\paragraph{Task.}
ERI separates method execution from admissibility:
\[
\begin{aligned}
\text{data} &\longrightarrow \text{diagnostic evidence}\\
&\longrightarrow \text{eligibility}
\longrightarrow \text{method or abstention}.
\end{aligned}
\]
Diagnostics record whether an assumption is supported, contradicted, or
inconclusive.  Eligibility rules cite this evidence and distinguish a primary
method from weakened or sensitivity-only alternatives.  If no admissible
primary method remains, or a fail-closed condition applies, the system
abstains.  Exact schemas, evidence fields, and abstention rules appear in
Supplement Sections A--B.

Typed evidence makes the decision boundary auditable without making the
infrastructure itself the scientific contribution.  Each item records the
claim tested, whether it is supported, contradicted, or inconclusive, the
diagnostic version, numerical result when available, and scope warnings.
Eligibility records
retain reasons and evidence references, while decisions distinguish primary,
weakened, and sensitivity-only procedures.  This representation lets the
evaluation locate an error in evidence acquisition, eligibility assessment,
or final selection.

The router is fail-closed.  Missing required evidence cannot be interpreted as
support, contradicted validity conditions cannot be overridden by a fluent
rationale, and inconclusive evidence remains distinct from both support and
contradiction.  Decisions preserve the versioned diagnostic result and the
evidence identifiers used by every eligibility judgment.  These constraints
make it possible to audit whether a failure arose before or after the
scientific boundary was evaluated.

\paragraph{Statistical setting.}
The task is inference for the unannualized sample Sharpe ratio of excess
returns.  Four registered candidates are IID Gaussian inference, Mertens
non-Gaussian variance with frequentist Probabilistic Sharpe Ratio, a
Bartlett/Newey--West HAC procedure, and circular-block bootstrap.  IID
Gaussian and Mertens procedures retain an IID basis; HAC and block bootstrap
permit serial dependence under their registered conditions.  Full definitions
and diagnostic thresholds appear in Supplement Sections B--C.  Diagnostics
cover input quality, dependence in returns and squared returns, distribution
shape, split-sample stability, and structured research provenance.  No HAC
bandwidth or bootstrap block length is inferred silently.

The dependence checks include returns and squared returns because conditional
heteroskedasticity may be visible in the latter even when linear
autocorrelation is weak.  HAC propagates both sources through the Sharpe
influence function.  Mertens uses plug-in skewness and raw kurtosis, while PSR
is interpreted as a frequentist probability under the registered
approximation rather than a posterior probability.  Bootstrap execution is
circular and requires an explicit block length, replication count, and seed.
These conventions are fixed before evaluation.

The frozen pilot uses two synthetic families and seeds 41100001--41100003.
\texttt{garch-persistent} exhibits persistent conditional
heteroskedasticity: IID procedures are ineligible, while HAC and block
bootstrap are eligible when the registered finite-fourth-moment condition
holds.  Under \texttt{structural-mean-break}, a stable full-sample Sharpe
target is not well defined, so all four procedures are ineligible and
abstention is required.  Hidden eligibility truth is never exposed to the
evaluated systems.

The structural-break family is deliberately different from a difficult but
estimable case.  When split-sample stability contradicts a common full-sample
target, robust standard errors do not restore the target's meaning.  The
correct benchmark action is therefore abstention, not escalation from IID to
HAC or bootstrap.  By contrast, persistent GARCH preserves a stable target but
requires dependence-robust methods.  Together, the families distinguish
method misspecification from target invalidity.

\paragraph{Information regimes.}
The 18 matched cells cross the two families and seeds with three information
regimes.  Full-information cells expose all structured diagnostics; budgeted
cells hide a dependence or stability group that can be acquired within a
fixed budget; open-world cells provide explicit, conflicting, or
instruction-like research provenance.  The complete cell registry and
visibility rules are in Supplement Section D.  These regimes test not only
whether a system recognizes a diagnostic result, but whether it requests
missing evidence within budget and resists treating provenance text as a new
instruction.

Budgeted observation makes evidence acquisition part of the decision while
keeping the allowed diagnostic set fixed.  Open-world provenance is data, not
an instruction channel: workflow prompts require the system to assess its
relevance and consistency without allowing embedded directive-like text to
change the task.  Full-information cells isolate routing after evidence is
already available.  Each scenario--seed combination appears in all three
regimes, yielding six cells per regime and 18 matched cells in total.

\paragraph{Evaluation.}
The benchmark loss \(L_i\) combines prespecified penalties for interval
behavior, invalid selection, harmful non-abstention, unnecessary abstention,
execution failure, and runtime.  A hidden-truth oracle selects the
minimum-loss eligible action for the matched cell.  Routing regret is
\[
\RoutingRegret_i=\max\{0,L_i-L_i^{\mathrm{oracle}}\},
\]
so lower is better.  An invalid primary selection chooses an ineligible
method, proceeds when the target is undefined, or fails execution; a claim is
valid only when execution succeeds and the primary selection is valid.  Loss
weights and oracle construction are policy choices documented in
Supplement Section D, not universal scientific constants.

All intended outputs remain in the denominator, including terminal failures.
The oracle is evaluator-only: it uses registered hidden truth and disjoint
simulation draws, never information available to the evaluated workflow.
Regret therefore measures excess decision loss relative to a matched-cell
reference, not whether a workflow can reproduce hidden labels from its visible
prompt.

The component penalties separate scientifically different errors.  Invalid
selection marks an ineligible primary method or continuation when the target
is undefined; harmful non-abstention identifies the latter failure
specifically; unnecessary abstention penalizes refusing when an eligible
primary method is available.  Claim validity requires successful execution and
a valid primary decision.  Runtime and interval behavior enter the registered
loss but do not replace these categorical validity outcomes.  Reporting both
loss and its components prevents a low average from masking an inadmissible
decision.
